%%%%%%%%%% source scan idx 51 | folio 44 %%%%%%%%%%

\underline{命题} 1.3. \underline{设 $f:X\to S$ 为固有光滑态射, $\underline W$ 为 $X$ 上秩为 $c$ 的局部自由层, $\underline F$ 为 $X$ 上的凝聚层, $d>0$ 为整数, $s:\underline W\to\underline O_X$ 为 $\underline W^\vee$ 的截面. 假设 $S$ 为诺特概形, 且满足:}

\begin{enumerate}[label=(\alph*),leftmargin=2.4em]
\item \underline{由 $s$ 的零点定义的 $X$ 的闭子概形 $H$ 在 $S$ 上光滑.}

\item \underline{在 $X$ 上局部地, 若 $s_i$ 为 $s$ 在 $\underline W^\vee$ 的一组基下的坐标, 则 $s_i$ 既为 $\underline O_X$-正则序列, 也为 $\underline F$-正则序列.}

\item \underline{对每个不全为零的整数序列 $k_\ell$, 都有}
\end{enumerate}

\begin{equation*}
R^if_*\left(\bigotimes_\ell\bigwedge^{k_\ell}\underline W
\otimes\Omega^j_{X/S}\otimes\underline F\right)=0
\qquad\underline{\text{当 }}i+j<d.
\end{equation*}

\underline{则:}

\begin{align*}
R^if_*\left(\Omega^j_{X/S}\otimes\underline F\right)
&\xrightarrow{\ \sim\ }
R^if_*\left(\Omega^j_{H/S}\otimes\underline F\right)
&&\underline{\text{当 }}i+j<d-c, \tag{i}\\
R^if_*\left(\Omega^j_{X/S}\otimes\underline F\right)
&\lhook\joinrel\longrightarrow
R^if_*\left(\Omega^j_{H/S}\otimes\underline F\right)
&&\underline{\text{当 }}i+j=d-c. \tag{ii}
\end{align*}

$H$ 在 $X$ 中的余法层 $N_{H/X}$ 同构于 $\underline W$ 在 $H$ 上的限制, 所以正合列

\begin{equation}\tag{1.3.1}
0\longrightarrow N_{H/X}\longrightarrow
\Omega^1_{X/S}\otimes\underline O_H\longrightarrow
\Omega^1_{H/S}\longrightarrow0
\end{equation}

在 $\Omega^j_{X/S}\otimes\underline O_H$ 上定义递减滤过, 其相继商为

\begin{equation}\tag{1.3.2}
\operatorname{Gr}^k\bigl(\Omega^j_{X/S}\otimes\underline O_H\bigr)
\simeq\bigwedge^k\underline W\otimes\Omega^{j-k}_{H/S}.
\end{equation}

由 (b) 还可得到

\begin{equation}\tag{1.3.3}
\operatorname{Gr}^k\bigl(\Omega^j_{X/S}\otimes\underline O_H
\otimes\underline F\bigr)
=\bigwedge^k\underline W\otimes\Omega^{j-k}_{H/S}
\otimes\underline F.
\end{equation}

我们对 $j$ 归纳, 同时对所有满足 (b) 和 (c) 的 $\underline F$ 证明 (i) 和 (ii). 当 $j<0$ 时, 两断言均为空断言.

%%%%%%%%%% source scan idx 52 | folio 45 %%%%%%%%%%

于是, 假定对于 $j'<j$ 已证明 (i) 和 (ii). 归纳假设适用于层 $\bigwedge^k\underline W\otimes\underline F$, 由 (c) 得到

\begin{equation}\tag{1.3.4}
R^if_*\left(\bigwedge^k\underline W\otimes
\Omega^{j-k}_{H/S}\otimes\underline F\right)=0
\quad\text{当 }k>0\text{ 且 }i+j<d-c+k.
\end{equation}

上面所考虑的 $\Omega^j_{X/S}\otimes\underline O_H\otimes\underline F$ 的滤过定义了谱序列

\begin{equation*}
E_1^{p,i-p}=R^if_*\left(\bigwedge^p\underline W\otimes
\Omega^{j-p}_{H/S}\otimes\underline F\right)
\Longrightarrow
R^if_*\left(\Omega^j_{X/S}\otimes\underline O_H
\otimes\underline F\right).
\end{equation*}

由 (1.3.4), 当 $p>0$ 且 $i+j\leq d-c$ 时 $E_1^{p,i-p}=0$; 因而

\begin{equation}\tag{1.3.5}
\left\{
\begin{aligned}
R^if_*\left(\Omega^j_{X/S}\otimes\underline O_H\otimes\underline F\right)
&\xrightarrow{\ \sim\ }
R^if_*\left(\Omega^j_{H/S}\otimes\underline F\right)
&&\text{若 }i+j<d-c,\\
R^if_*\left(\Omega^j_{X/S}\otimes\underline O_H\otimes\underline F\right)
&\lhook\joinrel\longrightarrow
R^if_*\left(\Omega^j_{H/S}\otimes\underline F\right)
&&\text{若 }i+j=d-c.
\end{aligned}
\right.
\end{equation}

由假设 (b), 赋予适当微分后, 层 $\bigwedge^k\underline W\otimes\Omega^j_{X/S}\otimes\underline F$ 构成 $\Omega^j_{X/S}\otimes\underline O_H\otimes\underline F$ 的左 Koszul 解消. 这个滤过定义了谱序列

\begin{equation}\tag{1.3.6}
E_1^{-p,q}=R^qf_*\left(\bigwedge^p\underline W\otimes
\Omega^j_{X/S}\otimes\underline F\right)
\Longrightarrow
R^{q-p}f_*\left(\Omega^j_{X/S}\otimes\underline O_H
\otimes\underline F\right).
\end{equation}

由假设 (c), 当 $p>0$ 且 $q+j<d$ 时 $E_1^{-p,q}=0$; 当 $p\notin[0,c]$ 时显然也为零. 因此,

\begin{equation}\tag{1.3.7}
\left\{
\begin{aligned}
R^if_*\left(\Omega^j_{X/S}\otimes\underline F\right)
&\xrightarrow{\ \sim\ }
R^if_*\left(\Omega^j_{X/S}\otimes\underline O_H\otimes\underline F\right)
&&\text{若 }i+j<d-c,\\
R^if_*\left(\Omega^j_{X/S}\otimes\underline F\right)
&\lhook\joinrel\longrightarrow
R^if_*\left(\Omega^j_{X/S}\otimes\underline O_H\otimes\underline F\right)
&&\text{若 }i+j=d-c.
\end{aligned}
\right.
\end{equation}
\EditorialFootnote{法文原印本与英文底本在 (1.3.7) 右侧均写作 $R^if_*(\Omega^j_{H/S}\otimes\underline F)$. 由 (1.3.6) 的收敛项, 此处应为 $R^if_*(\Omega^j_{X/S}\otimes\underline O_H\otimes\underline F)$; 再与 (1.3.5) 合成即得结论.}

现在由 (1.3.5) 和 (1.3.7) 即得断言 (i) 和 (ii).

